\chapter*{Conclusions}
\addcontentsline{toc}{chapter}{Conclusions}

This thesis mirrors the almost five years of this PhD.

Starting from a somehow canonical research question ("What the macroeconomic effects of a UBI are?") it delves first into a methodological journey and, later, in a modelling attempt to guess an answer.

Methodologically, it shifts from the naive enthusiasm and trust in quantitative models, to a deep scepticism and a movement towards more qualitative approaches.
This movement has been caused by the experience in attempting to model a complex enough economy on one side, and a broader feeling about empirical research and data quality and generalizability on the other.
On one hand, data are often too contextual to be used having the universality of the conclusion as the goal. On the other hand, models of an economy are either too simple (like the one presented in this model) or impossible to manage given their complexity (i.e. their non-linear and chaotic features) to be accurate tools (even if they are maybe the most accurate quantitative ones).

The model presented is thought of as a teaching tool for AB-SFC models, providing a somewhat realistic and simple model to start guessing the answer to the initial question.
However, if not constrained by the perimeter of what is researchable in economics today, I would argue that the initial question is less relevant that asking if an UBI is moral and desirable, and, in case of a positive answer, how we can collect the political power to implement it and observe \textit{in vivo} which are its effects.

In its simplicity the model is able to highlight that the answer to the original question is not unique, and it heavily depends on the other policies implemented alongside a UBI.
In fact the simple scenarios proposed diverge deeply: we observe both a scenario in which a UBI shifts the consumption towards the private sector, sustaining the aggregate demand but increasing the inequalities in wealth and income; and a scenario where a UBI implemented alongside other redistributive fiscal policies sustains the demand without increasing inequality.

To summarize, the methodological part of the thesis invites the researcher to recover qualitative analysis, at least alongside quantitative one, in order to better focus on which are the features of the system assumed to be relevant.
The second part tries to put on the ground this intuition to answer an economical question, providing in doing that a detailed example on how to implement an AB-SFC model that can be used as a teaching tool.
